% ============================================================
% Section 6: Analysis and Ablation
% ============================================================

\subsection{Component Contribution Analysis}
\label{sec:analysis:components}

The hierarchical system has three distinct architectural components beyond
the single-agent baseline: (1) multi-level domain decomposition, (2) the
Skeptical Senior Programmer verification loop, and (3) the reproduction
script as a machine-checkable verification contract. We analyze the
contribution of each via qualitative case analysis, as the 300-instance
evaluation protocol does not permit a full multi-arm ablation without
additional compute.

\paragraph{Benefit of domain decomposition.}
Consider an instance from \texttt{scikit-learn} (e.g.,
\texttt{scikit-learn\_\_scikit-learn-15512}) involving a bug in the
\texttt{RandomForestClassifier}'s warm-start validation logic. In the
baseline condition, the single agent receives the issue text and
pre-loaded file content in a single context and immediately attempts to
call \texttt{str\_replace}. If the pre-loaded file content is insufficient
(e.g., the relevant method is in a superclass not mentioned by name in the
issue), the agent produces an \texttt{old\_str} that does not match any
text in the file, causing \texttt{str\_replace} to fail, exhausting the
iteration budget on re-reads.

In the hierarchical condition, $\Dm_{\text{analysis}}$ runs
\texttt{code\_explorer} for up to 8 tool iterations---sufficient to trace
the inheritance chain across multiple files---and synthesizes a diagnostic
summary that precisely identifies the file and line range responsible. This
structured diagnostic is then passed to $\Dm_{\text{patch}}$, which
provides \texttt{patch\_writer} with a pre-loaded context window containing
only the relevant file section. The patch succeeds on the first attempt.

\paragraph{Benefit of the verification loop.}
In $k$ of the 53 successfully patched hierarchical instances, the first
patch attempt failed verification and a second attempt (with failure feedback
injected) succeeded. This retry behavior is not observable in the baseline
(which has no verification step). The structured feedback from
\texttt{test\_runner}---the exact error message from re-running
\texttt{/tmp/reproduce.py}---provides the patch writer with a precise signal
for correction that would otherwise require a fresh agent invocation.

\paragraph{Benefit of mandatory tool enforcement.}
In baseline runs, the model produces a text description of the fix (``I
would replace line 42 with...'') without calling \texttt{str\_replace} in
23.4\% of cases. With \texttt{mandatory\_tool\_call=True}, the tool loop
immediately aborts and returns an error prompt when the first response
contains no tool call, forcing a retry. This reduces text-only failure
responses to $<5\%$ in the hierarchical system's \texttt{patch\_writer}.

\subsection{Sympy Failure Analysis}
\label{sec:analysis:sympy}

All 77 sympy instances produce zero patches from either condition, with
sub-second elapsed times, indicating that the failure occurs before any LLM
call is made. Investigation reveals that all sympy instances fail at the
\texttt{clone\_and\_checkout} step: the base commits in SWE-bench Lite's
sympy instances pre-date the repository's shallow clone depth, making the
target commit unreachable without a full \texttt{git fetch --unshallow}
operation. The benchmark harness does not perform this operation, resulting
in a 100\% failure rate for sympy across both conditions.

This failure is repository-specific and systematic---it does not affect the
within-system comparison because both conditions fail identically on sympy
(contributing 77 tied pairs with no discordant outcomes). Nevertheless, it
represents a confound that should be addressed in future work by
implementing full-depth fetch fallback in the harness.

\subsection{Token Efficiency: Contextual vs.\ Total}
\label{sec:analysis:tokens}

We distinguish two distinct token efficiency metrics:

\paragraph{Total token footprint.} The hierarchical system generates
$55{,}558$ tokens on average versus $7{,}177$ for the baseline ($7.74\times$
higher). This overhead is the cost of running six agents with separate
planning and synthesis calls. For organizations with per-token cost
constraints, this represents a significant deployment consideration.

\paragraph{Per-agent contextual efficiency.} Each individual worker agent
operates with $3{,}388$ tokens on average, versus $7{,}177$ for the
baseline---a 52.8\% reduction. This is the operative quantity for reasoning
quality: as established in Section~\ref{sec:intro:failures}, reasoning
quality degrades as a function of total context length, not total tokens
generated across the pipeline.

The central efficiency insight is that the hierarchical system achieves
better per-agent context quality \emph{at the cost of} higher total token
consumption. This is an architectural tradeoff: pay more tokens overall to
ensure each individual reasoning step stays within the model's effective
context capacity.

\subsection{Qualitative Error Taxonomy}
\label{sec:analysis:errors}

We manually inspect 30 randomly sampled instances that both conditions fail
on (neither patches) and categorize the failure modes:

\begin{enumerate}
  \item \textbf{Repository setup failure} (sympy, 77 instances; 40.8\% of
    total failures): shallow clone prevents checkout. Excludes sympy from
    the categorical analysis below.

  \item \textbf{Code localization failure} (35\% of non-sympy failures):
    the relevant function is not in any file mentioned in the issue; neither
    condition navigates the inheritance/import chain deeply enough.

  \item \textbf{Semantic understanding failure} (28\% of non-sympy failures):
    the model correctly localizes the relevant file but misidentifies the
    root cause. Present in both conditions, but more common in the baseline
    (which lacks the \texttt{code\_explorer}'s iterative navigation).

  \item \textbf{Patch precision failure} (22\% of non-sympy failures):
    \texttt{str\_replace} receives an \texttt{old\_str} that does not
    match any text in the file (whitespace difference, version difference,
    or hallucinated line content). More common in the baseline condition
    (which relies solely on the pre-loaded content block).

  \item \textbf{Timeout} (15\% of non-sympy failures): the hierarchical
    system's per-instance timeout (600\,s) is exceeded on instances requiring
    large numbers of Domain Manager planning cycles.
\end{enumerate}
